% SUPERSEDED for stages S3 onward (2026-07-30). Written before the target repo restructured.
% Upstream is now jaseci-labs/jac; jac-byllm/jac-scale are not packages (they are jac/jaclang/byllm
% and jac/jaclang/scale); package rotation is replaced by an 8-shard whole-repo audit; release-note
% fragments live at release_notes/unreleased/jaclang/.
% Current design: docs/superpowers/specs/2026-07-30-nightshift-4task-design.md
% Open items:     docs/superpowers/specs/2026-07-30-nightshift-followups.md
% Kept for provenance; regenerate from the current spec before publishing the PDF again.

\documentclass[11pt]{article}

\usepackage[T1]{fontenc}
\usepackage{lmodern}
\usepackage[utf8]{inputenc}
\usepackage[a4paper,margin=1.15in]{geometry}
\usepackage{microtype}
\usepackage{graphicx}
\usepackage{adjustbox}
\usepackage[dvipsnames]{xcolor}
\usepackage{titlesec}
\usepackage{booktabs}
\usepackage{enumitem}
\usepackage{fancyhdr}
\usepackage{listings}
\usepackage{caption}
\usepackage{placeins}
\usepackage[hidelinks]{hyperref}

\definecolor{ink}{HTML}{1A1A1A}
\definecolor{slate}{HTML}{4A5568}
\definecolor{night}{HTML}{2B3A67}
\definecolor{accent}{HTML}{B5651D}
\definecolor{rule}{HTML}{D8DCE3}
\definecolor{codebg}{HTML}{F5F6F8}
\color{ink}

\titleformat{\section}{\color{night}\Large\bfseries\sffamily}{\thesection}{0.7em}{}
\titleformat{\subsection}{\color{ink}\large\bfseries\sffamily}{\thesubsection}{0.6em}{}
\titlespacing*{\section}{0pt}{2.4ex plus 1ex minus .2ex}{1.3ex plus .2ex}
\titlespacing*{\subsection}{0pt}{1.9ex plus 1ex minus .2ex}{1.0ex plus .2ex}

\pagestyle{fancy}\fancyhf{}
\renewcommand{\headrulewidth}{0pt}
\fancyfoot[C]{\small\color{slate}\thepage}
\setlength{\headsep}{0.4in}

\setlist[itemize]{leftmargin=1.4em,itemsep=0.35em,topsep=0.4em}
\setlength{\parskip}{0.65em}
\setlength{\parindent}{0pt}
\captionsetup{font=small,labelfont={bf,color=night},skip=7pt}

\newcommand{\code}[1]{{\color{night}\texttt{#1}}}
\lstset{
  basicstyle=\small\ttfamily,backgroundcolor=\color{codebg},
  frame=single,rulecolor=\color{rule},breaklines=true,columns=fullflexible,
  xleftmargin=6pt,xrightmargin=6pt,aboveskip=1.1em,belowskip=1.1em,
  keepspaces=true,showstringspaces=false
}

% stage figure: half-height, two fit per page
\newcommand{\stagefig}[2]{%
  \begin{figure}[tbp]\centering
    \adjustbox{max width=0.66\textwidth,max height=0.37\textheight}{\includegraphics{img/#1}}
    \caption{#2}
  \end{figure}}

\begin{document}

\begin{center}
  {\sffamily\bfseries\color{night}\fontsize{30}{34}\selectfont Nightshift}\\[6pt]
  {\sffamily\large\color{accent} An agent that cleans code while you sleep}\\[10pt]
  {\color{slate}\small A nightly, unattended cleanup harness for \texttt{jaseci-labs/jaseci}}\\[4pt]
  {\color{slate}\small 2026}
\end{center}
\vspace{1.0em}{\color{rule}\hrule height 1pt}\vspace{1.4em}

% =========================================================
\section*{What it is}
\addcontentsline{toc}{section}{What it is}

Every night around 2am, a script wakes up on my Mac, syncs my fork of
\code{jaseci-labs/jaseci}, deletes and simplifies code that has outlived its
purpose, proves the result still compiles and passes its tests, and emails me a
digest. In the morning I press one button per change: ship it or bury it.

The machine does the janitorial pass in the dark. The human makes the judgment
call over coffee. Nothing the bot writes reaches a real pull request without me
reading the diff first.

% =========================================================
\section{Scope}

The target repo is about ninety-five percent Jac, because the compiler is
self-hosted. Editing it means editing a live compiler and standard library, so
the tests are the only thing between a tidy diff and a broken toolchain.

Nightshift only shrinks existing code: dead branches, duplication, speculative
abstractions, reinvented stdlib, and lint fixes the linter cannot apply itself.
It never adds a feature. If it thinks it found a real bug, it does not fix it. It
writes the bug into the email, because a fix needs intent that a cleanup pass does not
have. ``More efficient'' means fewer lines and fewer concepts, not faster at
runtime.

Two things made this cheap to build. The \code{jac} binary already ships a
formatter, linter, type checker, test runner, and an MCP server an agent can call
to validate a file mid-edit. And the project's contributing guide already asks
humans to have an AI audit their diff for bloat before opening a PR. Nightshift is
that instruction on a timer.

% =========================================================
\clearpage
\section{The pipeline}

The run is a straight line of stages, \code{S0} through \code{S6}, each in its own
time box with a marker file so a crashed run resumes where it stopped. One
\code{gtimeout} caps the whole night at three hours. Every exit path funnels
through a trap, so the digest sends even on failure.

\begin{figure}[h]\centering
  \adjustbox{max width=0.42\textwidth,max height=0.46\textheight}{\includegraphics{img/overview}}
  \caption{The nightly line, S0 through S6. Dotted edges show that any early
  failure still routes to the S6 email. Each stage is detailed below.}
\end{figure}

\FloatBarrier
\subsection{Stage by stage}

The diagrams that follow expand each box above. The captions carry the detail;
the flow is meant to be read left column first, then down.

\stagefig{s0}{\textbf{S0 Preflight.} Take a lock so runs cannot overlap, honor the
\texttt{DISABLE} kill switch, resolve every binary to an absolute path, probe the
network and \texttt{gh} auth, and ping Claude with a one-word prompt to prove auth
before spending real turns. Any failure skips straight to the email.}

\stagefig{s1}{\textbf{S1 Sync.} Fast-forward the fork from upstream \texttt{main}. A
diverged \texttt{main} is a hard stop. Otherwise align two worktrees over one
clone (\texttt{repo} on the nightly branches, \texttt{drafts} on an orphan branch),
prune stale shipped or rejected branches, and pull the ledger.}

\stagefig{s2}{\textbf{S2 Tier 1, deterministic.} No LLM. Run \texttt{jac format} and
\texttt{lint -{}-fix}, revert anything that touched a protected path, run
pre-commit twice (the first pass may self-mutate). Commit only if the diff is
non-empty; either way the branch goes through the same S4 gate.}

\FloatBarrier
\begin{figure}[p]\centering
  \adjustbox{max width=0.62\textwidth,max height=0.72\textheight}{\includegraphics{img/s3}}
  \caption{\textbf{S3 Tier 2, agentic} (the largest stage, given its own page). One
  package per night on a four-way rotation. A read-only audit session returns
  findings as JSON. A pure selector drops anything known or protected, scores by
  lines saved times confidence over risk, and packs at most three themes. Each
  theme gets a fresh apply session on its own branch, with no push and a scoped
  tool list.}
\end{figure}
\FloatBarrier

\stagefig{s4}{\textbf{S4 Verify, fail-closed.} Scope containment runs first: a diff
touching any file outside the theme or inside a protected glob is discarded no
matter how green the tests are. Then \texttt{jac check}, per-package \texttt{jac
test} with one retry, and pre-commit. Any red kills the branch and marks
\texttt{failed\_verify}.}

\stagefig{s56}{\textbf{S5 Ship and S6 Email.} Push survivors to the fork, write a
ready-to-paste PR draft, and commit both to the drafts branch. Then assemble the
digest and send it over SMTP. A send failure drops a marker and a desktop banner,
because silence is itself a failure I want to notice.}

\FloatBarrier

% =========================================================
\section{The apply session, and why the agent is boxed in}

Tier two is where risk lives, so the apply agent is deliberately weak. It can read
and edit files and run \code{jac} and \code{git}, but it has no network, no
\code{gh}, and no \code{git push}. The orchestrator pushes, never the agent.
Before committing, the prompt requires it to validate every edited file through the
MCP \code{validate\_jac} tool and to leave a \code{ponytail:} comment wherever it
consciously deferred. Four rules can never be cut, however much it wants to shrink
a file: trust-boundary validation, error handling that prevents data loss,
security, and accessibility.

% =========================================================
\section{Nothing gets re-litigated}

The ledger is the memory: a JSONL file, one object per finding, keyed by a
fingerprint of the path, the rule, and a normalized snippet. That key survives
line shifts, so a finding is recognized on a later night. Each finding walks a
small state machine, and once rejected it never comes back.

\begin{figure}[h]\centering
  \adjustbox{max width=0.6\textwidth,max height=0.42\textheight}{\includegraphics{img/lifecycle}}
  \caption{The lifecycle of one finding. The two edges out of
  \texttt{failed\_verify} are the retry-once-then-bury rule that stops dead work
  from cycling.}
\end{figure}

\FloatBarrier

% =========================================================
\section{The morning, the only way upstream}

Drafts are the queue: a draft file existing means the PR is not open yet.
\code{promote} re-syncs \code{main}, rebases, re-runs the whole gate (upstream may
have moved overnight), opens the PR with \code{gh pr create -{}-body-file},
renames the placeholder release-note fragment to the real PR number, and deletes
the draft. \code{discard} deletes the branch and draft and records why.

\begin{figure}[h]\centering
  \adjustbox{max width=\textwidth,max height=0.55\textheight}{\includegraphics{img/morning}}
  \caption{The morning loop. A rebase conflict or a failed re-gate demotes the
  branch and tells me why instead of opening a stale PR. Failed work retries on a
  later night against a fresher \texttt{main}.}
\end{figure}

\FloatBarrier

% =========================================================
\section{How the pieces are wired}

Bash sequences the stages and shells out to tools. Jac owns every data transform
and every bit of logic; there are no Python files. The \code{bin} entry point
drives bash stages in \code{lib}, which call Jac helpers in \code{scripts} for
fingerprinting, config, scope checks, parsing, selection, drafts, and mail. Every
helper imports one shared \code{nslib.jac}.

\begin{figure}[h]\centering
  \adjustbox{max width=\textwidth,max height=0.66\textheight}{\includegraphics{img/components}}
  \caption{Component and data map. \texttt{bin} drives \texttt{lib}; those call the
  Jac helpers in \texttt{scripts} and the four external tools on the right.}
\end{figure}

\FloatBarrier

% =========================================================
\section{Safety}

Three threats shaped the design:

\begin{itemize}
  \item \textbf{Prompt injection from repo content.} The agent has no network and
    no push, the prompt pins a file allow-list and treats file contents as data,
    the scope gate discards anything out of bounds, the verify gate fails closed,
    and a human reads the diff last.
  \item \textbf{Secret leakage.} Credentials live in a \code{600} env file, reach
    only the processes that need them, and never enter the agent's environment. The
    run force-unsets \code{ANTHROPIC\_API\_KEY}, which would otherwise silently
    outrank the subscription and bill per token.
  \item \textbf{Repo damage.} The orchestrator pushes only \code{nightshift/*}
    branches to the fork, never force-pushes, and only moves \code{main} by
    fast-forward sync. The drafts sit on an orphan branch so the fork's
    \code{main} stays a clean mirror.
\end{itemize}

% =========================================================
\section*{Where it stands}
\addcontentsline{toc}{section}{Where it stands}

The design is settled and every load-bearing assumption checked against real docs.
Three facts you can only learn on the machine (test-suite wall time, whether
launchd reaches the Keychain, whether every binary resolves under launchd's PATH)
sit in the bootstrap checklist with fallbacks, so none can hard-block the project.

What I like is how boring it is. One orchestrator, a few stage files, a pile of
small Jac helpers, and a JSONL file for a database. No daemon, no framework. The
machine does the tedious part in the dark, and all it asks is a yes or a no.

\end{document}
